\section{An integral converse for the pure-coefficient curves}
\label{ci-section}

For unit residuals, every rational point above a positive finite base
coordinate has a genuine integral interpretation, including inputs
at the ramified prime or at infinity. Retain $\gamma=1+\zeta$,
$\gamma^2=3\zeta$, and integers $h,g\ge2$. Fix units $u_0,u_1$
and $e_0\in\{0,1\}$, and set
\[
 \tau_0=u_0\gamma^{e_0},\qquad\tau_1=u_1.
\]
Consider the full rational fiber product of Section~\ref{rd-section},
\begin{equation}\label{ci-fiber-product}
 L_{\tau_1,g}(t_1)=r(L_{\tau_0,h}(t_0)),\qquad
 r(x)=\frac{x}{x^2+x+1}.
\end{equation}

\begin{theorem}[Every positive point gives actual pure profiles]
\label{ci-pure-inverse}
For every rational point $(t_0,t_1)$ of \eqref{ci-fiber-product}
with $x=L_{\tau_0,h}(t_0)>0$ finite, write $x=a/b$ in lowest
positive terms. There are $e\in\{0,1\}$ and integers $R,Q\ge7$
such that
\begin{equation}\label{ci-pure-norms}
 a^2+ab+b^2=3^eR^h,\qquad F(a,b)=Q^g.
\end{equation}
Both are actual oriented profiles with unit residuals. No
unramifiedness or nonunit assumption is imposed on the input
coordinates. If $h$ is even, $e=e_0$; for odd $h$ the first
ramified parity may change.
\end{theorem}
\begin{proof}
Represent each rational projective input by a primitive integer
pair $(r_i,s_i)$ and put $w_i=r_i+s_i\zeta$. This includes
$s_i=0$ and the point at infinity. The element $w_i$ is nonzero.
Its $\gamma$-depth $\epsilon_i$ is zero or one: depth at least
two would imply that the rational prime three divides both
coordinates. Write
\[
 w_i=\gamma^{\epsilon_i}v_i,\qquad
 \epsilon_i\in\{0,1\},\quad v_i\text{ primitive unramified}.
\]
Dividing by $\gamma$ cannot introduce a common rational factor
into the coordinates, since that factor would then divide $w_i$.
The oriented UFD factorization shows that every power of $v_i$
is primitive and unramified.

For either output set $E=e_{\rm initial}+\epsilon_i n$, where
the initial exponent is $e_0$ on the first side and zero on the
second. The exact identity
\begin{equation}\label{ci-gamma-content}
 u\gamma^Ev_i^n
 =3^{\lfloor E/2\rfloor}u\zeta^{\lfloor E/2\rfloor}
                  \gamma^{E\bmod2}v_i^n
\end{equation}
has a primitive final factor: it has only one split orientation
at every prime and $\gamma$-depth at most one. Thus the full
content is exactly $3^{\lfloor E/2\rfloor}$, and the reduced norm is
\begin{equation}\label{ci-reduced-norm}
 3^{E\bmod2}(\operatorname N v_i)^n.
\end{equation}

Put $M=a^2+ab+b^2$. The first primitive output, after changing
the common sign if necessary, is exactly $(a,b)$. Therefore
\[
 M=3^eR^h,\qquad R=\operatorname N v_0,
 \qquad e=(e_0+\epsilon_0h)\bmod2.
\]
The second primitive output has ratio $ab/M$, with
$\gcd(ab,M)=1$, so its positive primitive coordinates are
$(ab,M)$. Consequently
\[
 F(a,b)=3^{\epsilon_1g\bmod2}Q^g,
 \qquad Q=\operatorname N v_1.
\]
Since $3\nmid F$, the displayed parity is zero. Also
$F(a,b)\ge13$, so $Q>1$. A nonunit primitive unramified
norm is an integer at least seven.

If $R=1$, then $M=3^e\le3$. Positive $a,b$ give $M\ge3$,
with equality only at $a=b=1$. This would imply $13=Q^g$,
impossible for $g\ge2$ and $Q>1$ because thirteen is prime.
Thus $R>1$ and $R\ge7$ as well. The primitive factorizations
remaining in \eqref{ci-gamma-content} give the actual oriented
identities after the unit sign adjustments. All projective input
cases have been included. The formula for $e$ proves the stated
parity distinction.
\end{proof}

Conversely, every actual positive pure-profile seed gives a
rational point of one of these curves by Sections~\ref{dc-section}
and~\ref{rd-section}. Hence, after allowing the finite unit choices
and both initial ramified parities, their positive rational points
correspond exactly to actual pure-profile seeds. This asserts
neither uniqueness of the point nor uniqueness of factorization.

\begin{corollary}[The even/even positive rational locus is empty]
\label{ci-even-even}
If both $h$ and $g$ are even, no curve \eqref{ci-fiber-product}
has a rational point whose base coordinate $x$ is positive and finite.
\end{corollary}
\begin{proof}
Such a point would give \eqref{ci-pure-norms}. Thus $MF/3^e$
would be a rational square. The exact combined-norm obstruction
of Theorem~\ref{qg-combined-obstruction} says that neither $MF$
nor $MF/3$ is a rational square for a positive primitive seed.
This is a contradiction.
\end{proof}

Only the positive rational base locus is excluded. Points over other
fields, negative coordinates and boundary base values are not ruled
out by this argument. The remaining exponent pairs stay open unless
a separate proved obstruction applies. The nonunit-multiplier
counterfamily of Section~\ref{pc-section} is consistent with this
converse: here the only cancellation is the ramified power of three
accounted for in \eqref{ci-gamma-content}. No general point-height
upper bound follows. These statements are ordinary results; the
complete geometric and oriented converse is not claimed as a Lean proof.
